\section{Introduction}

Variant calling is a cornerstone of bacterial genomics, supporting applications that range from outbreak investigation and identification of transmission clusters to antimicrobial resistance prediction and phylogenetic reconstruction \citep{didelot2012transforming,koser2012rapid,bush2020genomic}. In public-health and clinical laboratories, the confidence that clinicians and epidemiologists place in downstream conclusions---whether a neonatal MRSA cluster represents a single introduction, or whether an isolate carries a resistance allele that would change treatment---is bounded by the accuracy with which single-nucleotide polymorphisms (SNPs) and short insertions/deletions (indels) can be recovered from raw sequencing reads.

For the past decade and a half, short-read sequencing on Illumina instruments has been the mainstay of bacterial variant calling. Its relatively low per-base error rate and mature analysis ecosystem---exemplified by aligners such as \texttt{minimap2} \citep{li2018minimap2} and \texttt{BWA-MEM}, callers such as \texttt{BCFtools} \citep{danecek2021bcftools} and \texttt{FreeBayes} \citep{garrison2012freebayes}, and bacterial pipelines such as \texttt{Snippy} \citep{seemann_snippy}---have made Illumina the default reference point against which other technologies are measured. Large-scale benchmarking on Illumina data for bacterial genomes has clarified how genomic diversity and the choice of reference affect SNP-calling accuracy \citep{bush2020genomic}, setting clear expectations for what established short-read pipelines can and cannot do.

Oxford Nanopore Technologies (ONT) sequencing has emerged as an alternative with distinctive advantages for infectious-disease genomics: long reads that span repeats, near-real-time data streaming, and portable devices capable of running in remote field settings \citep{quick2016realtime,amarasinghe2020opportunities}. Historically, however, ONT base-calling accuracy lagged behind Illumina, limiting the technology's adoption for bacterial variant calling at clinical thresholds \citep{shafin2020nanopore}. The recent release of the R10.4.1 pore, together with the \texttt{Dorado} basecaller \citep{oxfordnanopore_dorado} that provides fast, high-accuracy (hac), and super-accuracy (sup) models, as well as duplex read calling that jointly decodes both strands of a single molecule, has narrowed the gap and in some cases exceeded Illumina per-read quality for bacterial genomes \citep{sereika2022oxford,wick2023perfect}.

In parallel, the variant-calling software landscape has been reshaped by deep-learning methods. \texttt{DeepVariant} demonstrated that a convolutional network over pileup images could match or surpass hand-crafted callers on Illumina data \citep{poplin2018deepvariant}. \texttt{Clair3} extended this idea to long reads by combining pileup and full-alignment networks \citep{zheng2022clair3}, and \texttt{NanoCaller} uses haplotype-aware deep networks tailored to long-read error profiles \citep{ahsan2021nanocaller}. ONT also distributes \texttt{Medaka} \citep{ont_medaka}, and \texttt{Longshot} remains a popular non-neural long-read caller that relies on haplotype-assisted SNV inference \citep{edge2019longshot}. Critically, all of these deep-learning callers ship with models trained exclusively on human sequencing data, leaving open the question of how well they generalise to bacteria, whose k-mer compositions, GC content, and DNA-modification patterns differ substantially from those of human genomes.

Despite active benchmarking of long-read variant callers on human genomes \citep{olson2022giab}, we lack a comprehensive, bacteria-focused evaluation of the latest ONT chemistry and software. Published bacterial benchmarks have almost exclusively considered Illumina data \citep{bush2020genomic}, and long-read benchmarks have concentrated on the human reference material curated by Genome in a Bottle \citep{olson2022giab}. In this study we fill that gap. We assess seven variant callers on ONT R10.4.1 data basecalled in fast, hac, and sup modes, in both simplex and duplex configurations, across 14 Gram-positive and Gram-negative bacterial species spanning a broad range of GC content. We compare against \texttt{Snippy} \citep{seemann_snippy} calls from paired deep Illumina sequencing. To avoid the limitations of fully simulated truthsets, we construct a pseudo-real benchmark by projecting variants from a closely related donor genome (average nucleotide identity near 99.5\%) onto each sample's gold-standard assembly, producing biologically realistic distributions of SNPs and short indels while still knowing the ground truth.

\section{Related Work}

Benchmarks of variant callers for long-read sequencing have largely centred on the human genome, leveraging reference materials such as HG002 from Genome in a Bottle \citep{olson2022giab}. These studies established the dominance of deep-learning methods for SNP calling but had limited power to speak to bacterial genomes, which differ in k-mer composition, repeat structure, and modification landscape. For bacterial variant calling specifically, Bush \emph{et al.} provided an extensive comparison of Illumina pipelines and showed that genomic diversity and the choice of reference have a larger effect on accuracy than the specific caller \citep{bush2020genomic}.

At the tool level, our study connects and contrasts several families of callers. Traditional probabilistic callers---\texttt{BCFtools/mpileup} \citep{danecek2021bcftools} and \texttt{FreeBayes} \citep{garrison2012freebayes}---remain widely deployed in bacterial pipelines, especially through wrappers such as \texttt{Snippy} \citep{seemann_snippy}. Long-read specific callers include \texttt{Longshot} \citep{edge2019longshot} for SNVs, and deep-learning callers \texttt{Clair3} \citep{zheng2022clair3}, \texttt{DeepVariant} \citep{poplin2018deepvariant}, \texttt{NanoCaller} \citep{ahsan2021nanocaller}, and \texttt{Medaka} \citep{ont_medaka}. Evaluation of these callers has been facilitated by recent advances in variant-comparison tools such as \texttt{vcfdist} \citep{dunn2024vcfdist}, which handle complex variants in a principled way.

Ground-truth construction for benchmarking is a non-trivial problem in itself \citep{olson2022giab}. Because calling variants against a sample's own reference yields none, benchmarks rely on either full simulation (which may not reflect biological distributions) or on gold-standard curated truthsets. Our pseudo-real approach---projecting a donor genome's real variants onto a polished reference---combines the control of simulation with the realism of biological variation, and follows in the spirit of prior bacterial benchmark constructions \citep{bush2020genomic}. It also depends on high-quality reference genomes, which we produce using the \texttt{Trycycler} consensus assembly workflow \citep{wick2021trycycler} followed by \texttt{Polypolish} short-read polishing \citep{wick2022polypolish} and Illumina-supported manual curation \citep{wick2023perfect}. Variant identification between reference and donor uses both \texttt{minimap2} \citep{li2018minimap2} and \texttt{MUMmer4} \citep{marcais2018mummer4}, with the intersection of the two callsets reducing aligner-specific artefacts. Genome quality is further screened with \texttt{CheckM} \citep{parks2015checkm}, reads are preprocessed with \texttt{fastp} \citep{chen2018fastp}, and repetitive regions are manipulated with \texttt{BEDTools} \citep{quinlan2010bedtools}. The resulting evaluation leverages the latest ONT R10.4.1 chemistry whose base-level accuracy has been characterised by Sereika \emph{et al.} \citep{sereika2022oxford}.
